\documentclass[11pt,a4paper]{article}
\usepackage[utf8]{inputenc}
\usepackage[portuguese-brazil]{babel}

% Matemática
\usepackage{amsmath, amssymb, amsthm}

% Tabelas e listas
\usepackage{booktabs}
\usepackage{array}
\usepackage{enumitem}

% Design
\usepackage{xcolor}
\usepackage[most]{tcolorbox}
\usepackage{geometry}
\geometry{margin=2.5cm}

% ============================================================
% Tema: troque estas três cores para mudar o visual do documento.
% ============================================================
\definecolor{accent}{HTML}{8B2B40}
\definecolor{accentDark}{HTML}{5D3A1A}
\definecolor{accentLight}{HTML}{F9F4E8}

% ============================================================
% Caixas de conteúdo: o menu "Caixas" da barra insere estes estilos.
% ============================================================
\tcbset{
    definitionBox/.style={
        enhanced, breakable, colback=accentLight, colframe=accent,
        fonttitle=\bfseries, coltitle=white, colbacktitle=accent,
        boxrule=1pt, arc=4pt,
        left=10pt, right=10pt, top=5pt, bottom=5pt,
    },
    noteBox/.style={
        enhanced, breakable, colback=white, colframe=accent,
        fonttitle=\bfseries, coltitle=accent, colbacktitle=white,
        boxrule=0pt, leftrule=4pt, sharp corners,
        left=10pt, right=10pt,
    },
    exampleBox/.style={
        enhanced, breakable, colback=accentLight!60!white, colframe=accentDark,
        fonttitle=\bfseries, coltitle=white,
        boxrule=0.5pt, arc=3pt,
        left=10pt, right=10pt,
    },
    textBox/.style={
        enhanced, breakable, colback=white, colframe=accentDark,
        boxrule=0pt, leftrule=2pt, sharp corners,
        left=10pt, right=10pt,
    },
    genericBox/.style={
        enhanced, breakable, colback=accentLight!60!white, colframe=accentDark,
        boxrule=0.5pt, arc=3pt,
        left=10pt, right=10pt,
    },
}

% Ambientes de teorema
\theoremstyle{plain}
\newtheorem{theorem}{Teorema}[section]
\theoremstyle{remark}
\newtheorem*{remark}{Comentário}

% Marcadores de lista coloridos
\setlist[itemize,1]{label={\color{accent}\textbullet}}
\setlist[enumerate,1]{label={\color{accent}\bfseries\arabic*.}}

% Bloco de título simples
\newcommand{\notesTitle}[2]{%
    \noindent{\color{accent}\rule{\linewidth}{3pt}}\par
    \vspace{1.2em}
    \begin{center}
        {\color{accentDark}\LARGE\bfseries #1}\par
        \vspace{0.5em}
        {\color{accent}\large #2}\par
        \vspace{0.4em}
        {\small \today}
    \end{center}
    \vspace{0.6em}
    \noindent{\color{accent}\rule{\linewidth}{1pt}}\par
    \vspace{1.5em}
}

% hyperref vem por último: ele altera vários outros pacotes.
\usepackage{xurl}
\usepackage[
    colorlinks=true,
    urlcolor=accent,
    linkcolor=accent,
    breaklinks=true,
    pdfpagemode=UseNone
]{hyperref}

\begin{document}

\notesTitle{{{title}}}{Álgebra Linear --- Autovalores e Autovetores}

\section{Motivação}

Muitas transformações lineares têm direções especiais que elas apenas \emph{esticam}, sem girar. Encontrar essas direções transforma um problema matricial difícil em alguns poucos problemas escalares. Pressione \textbf{Ctrl+Enter} para compilar e role qualquer um dos painéis: o outro acompanha, ancorado nos títulos das seções.

\begin{tcolorbox}[definitionBox, title=Definição 1.1]
Seja $A$ uma matriz $n \times n$. Um vetor não nulo $\mathbf{v}$ é um \emph{autovetor} de $A$ se
\[ A\mathbf{v} = \lambda \mathbf{v} \]
para algum escalar $\lambda$, chamado de \emph{autovalor} correspondente.
\end{tcolorbox}

\section{Encontrando autovalores}

Reescrevendo a definição como $(A - \lambda I)\mathbf{v} = \mathbf{0}$, vemos que existe solução não nula exatamente quando $A - \lambda I$ é singular.

\begin{theorem}
$\lambda$ é autovalor de $A$ se, e somente se,
\[ \det(A - \lambda I) = 0. \]
\end{theorem}

\begin{tcolorbox}[noteBox, title=Observação]
O polinômio $p(\lambda) = \det(A - \lambda I)$ é o \emph{polinômio característico}. Ele tem grau $n$, então $A$ tem no máximo $n$ autovalores distintos.
\end{tcolorbox}

\subsection{Um exemplo resolvido}

\begin{tcolorbox}[exampleBox, title=Exemplo 2.1]
Seja
\[ A = \begin{pmatrix} 2 & 1 \\ 1 & 2 \end{pmatrix}. \]
Então
\[ \det(A - \lambda I) = (2-\lambda)^2 - 1 = (\lambda - 1)(\lambda - 3), \]
logo os autovalores são $\lambda_1 = 1$ e $\lambda_2 = 3$, com autovetores $(1,-1)^{T}$ e $(1,1)^{T}$.
\end{tcolorbox}

\section{Resumo}

\begin{table}[h]
    \centering
    \begin{tabular}{@{}lll@{}}
        \toprule
        \textbf{Conceito} & \textbf{Equação} & \textbf{Significado} \\
        \midrule
        Autopar              & $A\mathbf{v} = \lambda\mathbf{v}$  & Direção só escalada \\
        Eq. característica   & $\det(A - \lambda I) = 0$          & Encontra todo $\lambda$ \\
        Traço                & $\operatorname{tr} A = \sum \lambda_i$ & Soma dos autovalores \\
        Determinante         & $\det A = \prod \lambda_i$         & Produto dos autovalores \\
        \bottomrule
    \end{tabular}
\end{table}

\begin{remark}
Matrizes reais simétricas, como a do Exemplo 2.1, sempre têm autovalores reais e autovetores ortogonais.
\end{remark}

\section{Exercícios}

\begin{enumerate}
    \item Encontre os autovalores de $\begin{pmatrix} 4 & 0 \\ 0 & -1 \end{pmatrix}$.
    \item Qual das opções é o polinômio característico de $\begin{pmatrix} 1 & 2 \\ 0 & 3 \end{pmatrix}$?
    \begin{enumerate}[label=\alph*),leftmargin=1.5em,topsep=0.3em,itemsep=0.15em]
        \item $\lambda^2 - 4\lambda + 3$
        \item $\lambda^2 + 4\lambda + 3$
        \item $\lambda^2 - 3$
    \end{enumerate}
\end{enumerate}

\vspace{1em}
\noindent Leitura complementar: \href{https://pt.wikipedia.org/wiki/Autovalores_e_autovetores}{Autovalores e autovetores}.

\end{document}
